% Master's thesis -- UMONS, Faculté des Sciences
% Based on the UMONS master's thesis template (Memoire-canevas, A. Buys).
\documentclass[a4paper,11pt,notitlepage]{report}

\usepackage[utf8]{inputenc}
\usepackage[T1]{fontenc}
\usepackage{amsmath}
\usepackage{amsthm}
\usepackage{amssymb}
\usepackage{enumerate}
\usepackage{graphicx}
\usepackage[vmargin=2.6cm,hmargin=2.9cm]{geometry}
\usepackage{indentfirst}
\usepackage{cite}
\usepackage[english]{babel}
\usepackage{textcomp}

\usepackage{pslatex}

\usepackage{xcolor}
\usepackage{listings}
\usepackage{url}
\usepackage[hidelinks]{hyperref}

% ---------------------------------------------------------------------------
% YAML listings
% ---------------------------------------------------------------------------
\lstdefinelanguage{yaml}{
  keywords={true,false,null},
  comment=[l]{\#},
  morestring=[b]', morestring=[b]",
  sensitive=true
}
\lstset{
  language=yaml,
  basicstyle=\ttfamily\small,
  columns=fullflexible,
  keepspaces=true,
  frame=single, rulecolor=\color{gray!50},
  commentstyle=\color{gray}\itshape,
  stringstyle=\color{teal!70!black},
  breaklines=true,
  showstringspaces=false
}

\pagestyle{myheadings}

\newcommand{\STitre}{UNIVERSIT\'E DE MONS}
\newcommand{\FTitre}{FACULT\'E DES SCIENCES}
\newcommand{\Titre}{Refactoring GitHub Actions for Reuse}

\title{
  \begin{minipage}\linewidth\centering
  \LARGE\bf\STitre
  \vskip20pt
  \Large\bf\FTitre
  \vskip20pt
  \Huge \Titre
  \end{minipage}
}

\author{Abdellah BOUMAIZA\\[1.5ex]
Master en Sciences Informatiques\\[1.5ex]
{\tt \small abdellah.boumaiza@student.umons.ac.be}
}
\date{Academic year 2026--2027}
\begin{document}

\maketitle
\smallskip

\markright{\Titre}

\thispagestyle{empty}

\bigskip

\vfill

\begin{quote}
{Master's thesis carried out in the {\bf Software Engineering Lab}, as part of the
{\bf Master en Sciences Informatiques}. \\ }
\end{quote}
\begin{quote}
{\centering \bf Supervisor: } Prof.\ Tom MENS \\
{\centering \bf Co-supervisor: } Aref TALEBZADEH BARDSIRI
\end{quote}

\vfill

\includegraphics[width=0.2\textwidth]{img/FS.pdf} \hspace{0.6\textwidth} \includegraphics[width=0.2\textwidth]{img/logo_UMONS.pdf} \hspace{0.2\textwidth}

\newpage
The author, BOUMAIZA Abdellah, declares having complied with the ethical rules in force,
including the University's charter on the use of Artificial Intelligence.

% cf. thesis regulations

\newpage
\tableofcontents
\normalsize
\newpage

\setlength{\parindent}{12pt}
\setlength{\parskip}{1ex plus 0.5ex minus 0.2ex}
\setlength{\emergencystretch}{1em}

\addcontentsline{toc}{chapter}{Introduction}

\section*{Introduction} \label{sec_intro}

This thesis aims to design and implement an automated refactoring tool for GitHub
Actions workflows that introduces the platform's reuse mechanisms, namely composite
actions and reusable workflows, while preserving workflow behaviour. The scope is
initially restricted to reuse-oriented refactorings.

% ===========================================================================
\chapter{Reuse mechanisms in GitHub Actions} \label{ch_reuse}
% ===========================================================================

Before designing refactorings, it is necessary to establish which reuse mechanisms the
platform already offers, since each of them defines a possible target for an automated
transformation.

\begin{itemize}
  \item \textbf{Actions} encapsulate a single reusable unit of logic, referenced in a step
    through \texttt{uses:\ owner/repo@ref}. They are widely adopted, mostly through the
    GitHub Marketplace.
  \item \textbf{Composite actions} group a sequence of steps in an \texttt{action.yml} file
    and execute within the calling job. They accept inputs and expose outputs, but cannot
    access the \texttt{secrets} context and require an explicit \texttt{shell} for every
    \texttt{run} step. They can be stored locally in the repository or in a separate one.
  \item \textbf{Reusable workflows}, declared with the \texttt{workflow\_call} trigger, are
    invoked at the job level and therefore encapsulate entire jobs. They accept inputs,
    secrets and outputs, but secrets must be passed explicitly or through
    \texttt{secrets:\ inherit}.
  \item \textbf{Matrix strategies} execute a single job definition over several parameter
    values, which makes them suitable for replacing near-identical jobs within the same
    workflow.
\end{itemize}

These mechanisms operate at different levels of granularity, namely the step, the job and
the workflow. This granularity determines which kind of duplication each of them can
remove, and consequently which refactoring should be applied to a given clone.

% ===========================================================================
\chapter{Related work and its relevance} \label{ch_related}
% ===========================================================================

\section{Code clones in workflows} \label{sec_clones}

Cardoen, Mens and Decan~\cite{Cardoen2026} analysed 118K workflows and found clones in
67\% of them, most of which are short sequences of two or three steps. They estimate that
local composite actions could remove 83.7\% of clones at file scope and 75.6\% at
repository scope, whereas reusable workflows would remove only about 2\%, since they
require entire jobs.
\emph{Relevance:} this study identifies \emph{Extract Composite Action} as the refactoring
with the highest potential and provides both refactoring candidates and a quantitative
baseline. Its estimate, however, does not verify whether the reuse would preserve
behaviour, which is precisely the gap this thesis addresses.

\section{Practitioners' reuse practices} \label{sec_practices}

Onsori Delicheh et al.~\cite{OnsoriDelicheh2026} surveyed 419 practitioners. Copy-paste
remains common for reasons of convenience and control; 26.9\% of respondents were unaware
of composite actions and 25.4\% unfamiliar with them, and reuse is often perceived as
adding complexity.
\emph{Relevance:} the main barriers are awareness and effort rather than trust, which a
tool generating reuse automatically can directly lower, provided that its output remains
readable and that it suggests changes rather than imposing them.

\section{Refactoring and workflow smells} \label{sec_smells}

Fowler~\cite{Fowler2018} defines refactoring as a behaviour-preserving improvement of
internal structure and catalogues refactorings together with the smells they remove;
\emph{Extract Function}, for instance, directly corresponds to \emph{Extract Composite
Action}. Khatami et al.~\cite{Khatami2024} propose a catalogue of smells specific to
GitHub Actions workflows.
\emph{Relevance:} these works provide the conceptual framework of the thesis, namely the
definition of a refactoring, the link between smells and refactorings, and the
duplication-related smells that the tool should target.

% ===========================================================================
\chapter{Existing tools} \label{ch_tools}
% ===========================================================================

A workflow refactoring tool requires three capabilities: the faithful parsing and
serialisation of YAML, the querying and transformation of its structure, and an
understanding of workflow semantics.

\section{YAML processing} \label{sec_yaml}

\textbf{PyYAML} is the reference YAML library for Python. It maps documents onto native
dictionaries and lists, which facilitates analysis, but discards comments, key order and
formatting upon serialisation. Moreover, it implements YAML~1.1, under which the
\texttt{on} key of a workflow is interpreted as the boolean \texttt{True}.
\textbf{ruamel.yaml}, in contrast, provides a round-trip mode in which documents are
loaded into dedicated structures that retain comments, key order, quoting style and
anchors. Its output thus remains close to the original, which is essential for developer
acceptance and makes it the natural foundation for the tool.

\section{Querying and transforming YAML} \label{sec_query}

Whereas XML benefits from XPath and XQuery for querying and XSLT for transformation, YAML
lacks standard equivalents. \textbf{JSONPath} and \textbf{JMESPath} support declarative
queries over YAML documents once loaded as data, but provide no transformation
capabilities. \textbf{yq} supports both querying and editing from the command line, yet
offers limited control over formatting and has no notion of workflow semantics.
A promising direction is therefore a lightweight query and rewriting layer, built in
Python on ruamel.yaml and tailored to GitHub Actions workflows.

\section{Workflow analysis tools} \label{sec_analysis}

\textbf{gawd}~\cite{RostamiMazrae2024}, developed within the Software Engineering Lab,
represents workflows as trees and identifies the items added, removed, moved, renamed or
modified between two versions. It may inform the tool's internal representation and help
verify that a refactoring yields exactly the intended changes.
\textbf{actionlint} is a static checker that verifies workflow syntax, the types of
\texttt{\$\{\{~\}\}} expressions, job dependencies, action inputs and embedded shell
scripts. As its checks target errors rather than design issues such as duplication, it is
better suited to validating refactored workflows.
Finally, the \textbf{clone detector of Cardoen et al.}~\cite{Cardoen2026} identifies
identical and near-identical fragments by hashing normalised workflow trees, thereby
providing both a source of refactoring candidates and a baseline for evaluation.

% ===========================================================================
\chapter{Research direction} \label{ch_direction}
% ===========================================================================

\section{Candidate refactorings} \label{sec_candidates}

Following Fowler~\cite{Fowler2018}, a refactoring is defined as a transformation that
improves the internal structure of a program without altering its observable behaviour.
For GitHub Actions, this suggests several reuse-oriented refactorings: the extraction of
duplicated steps into a local composite action (and its inverse, inlining), the extraction
of duplicated jobs into a reusable workflow invoked through \texttt{workflow\_call}, the
splitting and merging of jobs and steps, and the decomposition of a complex workflow into
two interdependent workflows. Each refactoring will be characterised by the duplication or
smell it addresses, its preconditions, and the transformation it performs. Given that most
clones consist of short step sequences, \emph{Extract Composite Action} is the natural
starting point.

\section{Technical approach} \label{sec_approach}

Following the classical architecture of refactoring tools, each workflow is parsed with
ruamel.yaml into a tree preserving its concrete syntax. This tree is then abstracted into a
workflow model capturing jobs, steps, the dependency graph induced by \texttt{needs}, and
the references contained in expressions. Refactoring opportunities are identified through
queries on this model, the corresponding transformations are applied, and the result is
serialised back to YAML together with any newly created file, such as an
\texttt{action.yml}. The envisaged outcome is a Python command-line tool, comparable in
design to gawd, capable of suggesting or applying refactorings.

\section{Behaviour preservation} \label{sec_preservation}

The principal challenge lies in guaranteeing that a refactored workflow remains
behaviourally equivalent to the original. In GitHub Actions, this is non-trivial: each job
executes on a separate runner and therefore does not share its file system; reusable
workflows do not inherit secrets unless these are passed explicitly; and composite actions
require a \texttt{shell} attribute on every \texttt{run} step and cannot access the
\texttt{secrets} context directly. Such constraints will be formalised as preconditions to
be verified before any refactoring is applied, and refactored workflows will additionally
be validated with actionlint.

For instance, consider a duplicated step that reads a secret: \texttt{run:\ ./deploy.sh}
with \texttt{env:\ \{ TOKEN:\ \$\{\{ secrets.DEPLOY\_TOKEN \}\} \}}. Moved as is into a
composite action, it would fail, since the \texttt{secrets} context is unavailable there.
A behaviour-preserving extraction must instead turn the secret into an input of the action
and pass it explicitly from the caller:

\begin{lstlisting}
# .github/actions/deploy/action.yml
inputs:
  token: { required: true }
runs:
  using: composite
  steps:
    - run: ./deploy.sh
      shell: bash
      env: { TOKEN: "${{ inputs.token }}" }

# calling job
- uses: ./.github/actions/deploy
  with: { token: "${{ secrets.DEPLOY_TOKEN }}" }
\end{lstlisting}


\clearpage
\addcontentsline{toc}{chapter}{Bibliography}
\bibliographystyle{acm}
\bibliography{references}

\end{document}
